\section{Experiments}
\label{sec:experiments}

We evaluate H2O+ on a SUMO microsimulation of a real-topology 12-line bus corridor with five questions in mind.
\textbf{Q1: Headline performance.}
At the default operating point, how does H2O+ compare against pure offline RL, pure online RL, and the RE-SAC SUMO expert checkpoint (ep39) on cumulative cost?
\textbf{Q2: Component recipe.}
Among the H2O+ component choices (which discriminator, whether to add the offline-Q floor, KL anchor, warmup), which combinations help in the transit MDP and which do not---and does the answer depend on whether Stage~2 trains in the same simulator used for evaluation (oracle) or in a lower-fidelity simulator (deployment-relevant)?
\textbf{Q3: Multi-scenario robustness and operational metrics.}
Does the headline ranking hold across SUMO seeds and OD-scale perturbations under common random numbers, and how do the methods compare on passenger-side metrics (waiting time, headway CV, Jain fairness)?
\textbf{Q4: Assumption 1 data-level check.}
Does the action-conditional domain invariance assumption hold empirically on the benchmark, verified \emph{from data} rather than from the learned discriminator's own output?
\textbf{Q5: Sensitivity and generality.}
Are the main results stable to offline-data composition, ensemble size, KL anchor, and offline-data scale?

\subsection{Benchmark setup}
\label{sec:exp:setup}

\paragraph{Evaluation environment}
We use SUMO 1.15 \citep{lopez2018microscopic} with \texttt{LIBSUMO\_\allowbreak AS\_\allowbreak TRACI=1} to simulate a real 12-line corridor with 389 vehicles and a fixed passenger OD profile over an 18\,000-second operating period.
The network includes 7X-series routes with their actual station sequences, headway-scheduled dispatch, and calibrated dwell dynamics.
Numbers in \Cref{tab:main} (default operating point) use a single SUMO seed and the default OD profile for reproducibility; \Cref{tab:multiseed} (\Cref{sec:exp:multiseed}) reports the same methods averaged across $3$ SUMO seeds $\times$ $3$ OD-scale multipliers under common random numbers.

\paragraph{Low-fidelity simulator}
The sim-core environment runs an LSTM travel-time model with analytic boarding/alighting; at the same step length, it runs approximately $40\times$ faster than SUMO per event.
As reported in \Cref{sec:problem:h2o}, its reward magnitude is systematically smaller than SUMO's: $\lvert\E_{\Msim}[r]\rvert / \lvert\E_{\Mreal}[r]\rvert \approx 0.2$ under matched initial conditions.

\paragraph{Offline dataset}
$\Doff$ contains $675{,}000$ transitions from four behaviour policies:
SAC-trained (180K), headway-threshold heuristic (160K), $\epsilon$-random (180K), and the best RE-SAC checkpoint from a prior SUMO-environment run, which we refer to as \emph{ep39} (155K).
We note explicitly that ep39 transitions are part of $\Doff$, so the offline-RL and H2O+ methods we report below have observed ep39 behaviour at training time; this means that improvements over ep39 should be read as policy improvement after distilling a strong RE-SAC expert together with other behaviour data, not as discovering a strategy independent of that expert.
The composition is intentionally heterogeneous to allow Q-floor anchoring and Cal-QL pessimism to interact with multi-policy support.

\paragraph{Metrics}
Beyond cumulative episode reward, the evaluation harness \texttt{eval\_\allowbreak with\_\allowbreak metrics.py} produces a fuller set of operational indicators we report selectively:
per-decision cost (reward divided by number of decisions),
per-leg passenger waiting time (boarding minus stop-arrival, excluding transfer),
per-line headway coefficient of variation (HW CV),
\emph{bunching rate} (fraction of station departures with $h_{\text{prev}} < 180$\,s),
\emph{severe bunching rate} ($h_{\text{prev}} < 90$\,s),
\emph{large-gap rate} (fraction of headways exceeding $1.5\times$ scheduled),
and Jain's fairness index over per-line HW CV.
Tab.~\ref{tab:main} reports cumulative reward and per-decision cost (single-scenario summary).
\Cref{tab:multiseed} reports cumulative reward, passenger waiting, HW CV, large-gap rate, and Jain (multi-scenario operational comparison); bunching/severe-bunching rates are dominated by the same headway-CV signal and are reported for the SUMO-online ablation only (\Cref{tab:ablation}).

\paragraph{Training protocol}
Each configuration trains for 200 epochs with 100 SAC updates and 100 online rollout events per epoch.
An $\Ecnt{=}5$ ensemble is used throughout unless otherwise specified; pessimism coefficient $\beta_{\text{LCB}}{=}1.0$; IS-weight clipping $[0.1, 10]$; Q-floor weight $\alpha_{\text{cal}}{=}5.0$; SAC learning rates $3\mathrm{e}{-4}$.
Every configuration is run with five random seeds ($\{42, 123, 456, 789, 2024\}$) and we report mean $\pm$ standard error, with pairwise Wilcoxon signed-rank tests for method comparison.
Complete hyperparameters appear in \Cref{app:impl}.

\subsection{Baselines}
\label{sec:exp:baselines}

The methods reported in \Cref{tab:main} fall into four categories:
(i) \emph{rule-based / no-learning and expert reference}: zero-hold (no intervention) and the RE-SAC SUMO expert checkpoint ep39 (a learned policy, but included as a fixed deterministic checkpoint without further training; its transitions are also part of $\Doff$).
(ii) \emph{pure online RL}: SAC trained from scratch in the cheap simulator $\Msim$ for 200 epochs with no offline data and no pretrained checkpoint.
(iii) \emph{pure offline RL}: Cal-QL on the full $\Doff$ buffer.
(iv) \emph{H2O+ family}: hybrid offline-pretrain plus Stage~2 online interaction, parameterised by which discriminator and which optional anchors are used.
Within (iv) we report two settings of the Stage~2 online environment (SIM-online for the deployment-relevant case; SUMO-online as a same-fidelity oracle); the discriminator and anchor ablation is detailed in \Cref{sec:exp:ablations}.

Behaviour cloning on the full $\Doff$, behaviour cloning on the ep39-only subset of $\Doff$ ($3$ seeds), the deterministic RE-SAC SUMO expert ep39 checkpoint, and Daganzo cooperative holding at $\alpha \in \{0.3, 0.4, 0.6\}$ \citep{daganzo2009headway, xuan2011dynamic} are included in \Cref{tab:multiseed} (\Cref{sec:exp:multiseed}).
We additionally train and evaluate five standard offline-to-online RL baselines on the same $\Doff$ buffer with three training seeds each: IQL \citep{kostrikov2021offline}, AWAC \citep{nair2020awac}, TD3+BC \citep{fujimoto2021minimalist} (pure-offline), and RLPD \citep{ball2023efficient}, WSRL \citep{zhou2024wsrl} (offline+online hybrid). The implementations and results are reported in \Cref{tab:multiseed,tab:deployment_paired,tab:passenger_full} and discussed in \Cref{sec:exp:multiseed}; collectively, they directly address the question of whether H2O+'s dynamics-correction is necessary or whether simpler hybrids suffice.

\subsection{Main results}
\label{sec:exp:main}

\Cref{tab:main} compares H2O+ against four reference classes on the high-fidelity SUMO target (18\,000 s episode).
The deployment-relevant H2O+ setting is \emph{SIM-online}: Stage~2 trains in the cheap LSTM sim-core $\Msim$ (matching the assumption that the operator cannot put a learning agent on the live fleet or in a high-fidelity simulator continuously); evaluation is in $\Mreal$.
A second \emph{SUMO-online} group is reported as an upper-bound oracle in which the policy is trained directly in the same high-fidelity environment used for evaluation, indicating what could be reached if compute were unconstrained.

The key empirical claims:
(i) \emph{pure online RL}---SAC trained from scratch in $\Msim$ with no offline data---fails on this benchmark, $-1{,}654$K mean across 5 seeds, with one seed diverging to $-2{,}309$K (worse than zero-hold).
(ii) \emph{pure offline RL} (Cal-QL on the full $675$K-transition buffer) reaches $-749 \pm 26$K mean over 5 seeds: $\approx 2.2\times$ better than pure online, but still below the RE-SAC SUMO expert ep39 ($-666$K, single deterministic episode).
(iii) Under the deployment-relevant SIM-online setting, the best H2O+ configuration (Contrastive InfoNCE discriminator without Q-floor) reaches $-658$K mean over 5 seeds, slightly exceeding the RE-SAC expert in this default-scenario protocol and improving over pure offline by $\approx 12\%$; this single-scenario ordering does not generalise to the multi-scenario stress grid (\Cref{tab:multiseed}), where ep39 leads under low/moderate demand and H2O+ leads under high demand.
(iv) Under the SUMO-online oracle setting, H2O+ with the simple TransitionDiscriminator reaches $-646 \pm 16$K mean over 4 seeds, an additional $\approx 2\%$ improvement that quantifies the gap due to using the cheap simulator instead of the target.

\begin{table*}[t]
\centering
\small
\caption{Main results on $\Mreal$ (high-fidelity SUMO target, 18\,000 s episode, single SUMO seed = $42$, default OD profile). Cumulative reward, higher is better. Boldface marks the best deployment-relevant (SIM-online) and oracle (SUMO-online) results respectively. The reported $\pm$ ranges reflect training-seed variance only (4--5 training seeds depending on row; the discrepancy reflects which configurations completed all 5 seeds vs lost one to early job termination, not a different evaluation protocol).}
\label{tab:main}
\resizebox{\textwidth}{!}{%
\begin{tabular}{lrrr}
\toprule
Method & Cum.\ reward (K) & Per-step & Seeds \\
\midrule
\multicolumn{4}{l}{\emph{Reference}} \\
Zero-hold (deterministic)                            & $-1{,}682$            & $-136$         & --- \\
RE-SAC SUMO expert ep39 (single deterministic run)   & $-666$                & $-54$          & --- \\
\midrule
\multicolumn{4}{l}{\emph{Pure offline / pure online}} \\
Pure online RL (SAC scratch in $\Msim$)              & $-1{,}654 \pm 329$    & $-138$         & 5 \\
Pure offline RL (Cal-QL on $\Doff$, $675$K)           & $-749 \pm 26$         & $-61$          & 5 \\
\midrule
\multicolumn{4}{l}{\emph{H2O+ (SIM-online: deployment-relevant)}} \\
Contrastive only                                     & $\mathbf{-658 \pm 23}$ & $\mathbf{-54}$ & 5 \\
DynamicsDisc (temp $0.5$)                            & $-669 \pm 25$          & $-55$          & 5 \\
TransitionDisc only                                  & $-674 \pm 28$          & $-55$          & 4 \\
Contrastive + Q-floor + KL (full)                    & $-705 \pm 38$          & $-58$          & 5 \\
\midrule
\multicolumn{4}{l}{\emph{H2O+ (SUMO-online: same-fidelity oracle, not deployable)}} \\
TransitionDisc only (oracle best)                    & $\mathbf{-646 \pm 16}$ & $\mathbf{-53}$ & 4 \\
Contrastive + Q-floor + KL (full)                    & $-663 \pm 16$          & $-54$          & 4 \\
TransitionDisc + Q-floor                             & $-676 \pm 60$          & $-55$          & 4 \\
\bottomrule
\end{tabular}%
}
\end{table*}

\subsection{Ablation study}
\label{sec:exp:ablations}

We run an ablation sweep of approximately 30 distinct configurations across five groups (3 baselines, 9 SUMO-online, 9 SIM-online, 3 source-removal, 4--6 data-size points), each with 4--5 seeds, totalling on the order of 150 training runs, to test each component, the simulator-fidelity transfer claim, and the data-efficiency story.
Group A (offline-only) establishes the offline ceiling under three offline-RL baselines.
Groups B and C vary the online environment (SUMO vs SIM) while toggling the discriminator choice (none / TransitionDisc / DynamicsDisc / Contrastive), the offline-Q floor, the KL anchor, and (Group B only) sim-reward rescaling.
Group D removes one offline-data \emph{source} at a time (SAC-only, heuristic-only, weak-only) to probe dataset composition robustness.
Group E varies the offline-data \emph{size} at four levels (100K, 200K, 400K, full $\approx$ 675K transitions) under both pure offline (Cal-QL) and the full H2O+ pipeline, to test whether the online phase compensates for offline-data scarcity.

\Cref{tab:ablation} reports the nine SUMO-online configurations averaged across 4 seeds.
Three observations stand out.
First, the simplest H2O+ variant---TransitionDiscriminator IS reweighting with no Q-floor---is the strongest single configuration ($-646$K), narrowly beating all combinations that add Q-floor or replace the discriminator.
Second, the offline-Q floor's contribution is configuration-dependent: it improves \texttt{sumo\_baseline} ($-710 \to -678$K, $+32$K) and \texttt{sumo\_is} ($-694 \to -676$K, $+18$K) but \emph{harms} \texttt{sumo\_darc} ($-646 \to -676$K, $-30$K), so adding Q-floor has no consistent positive effect across IS-weight regimes.
Third, the factored DARC variant (DynamicsDiscriminator) is the worst configuration ($-720$K), indicating that the explicit action-conditional factorisation underperforms the simpler binary classifier in this domain.

\begin{table*}[t]
\centering
\small
\caption{Ablation on SUMO online (Group B), mean cumulative reward over 4 seeds (42, 123, 456, 789). Bold: best configuration. Rows sorted by mean reward.}
\label{tab:ablation}
\resizebox{\textwidth}{!}{%
\begin{tabular}{llr}
\toprule
Config & Components & Cum.\ reward (K) \\
\midrule
\texttt{sumo\_darc}         & TransDisc IS                       & $\mathbf{-646}$ \\
\texttt{sumo\_full}         & Contrastive + Q-floor + KL + warmup & $-663$ \\
\texttt{sumo\_darc\_calql}  & TransDisc + Q-floor                & $-676$ \\
\texttt{sumo\_is\_calql}    & Contrastive + Q-floor              & $-676$ \\
\texttt{sumo\_calql}        & no IS, Q-floor only                & $-678$ \\
\texttt{sumo\_is}           & Contrastive only                   & $-694$ \\
\texttt{sumo\_rescale}      & TransDisc + sim-reward rescale     & $-702$ \\
\texttt{sumo\_baseline}     & no IS, no Q-floor                  & $-710$ \\
\texttt{sumo\_dyn}          & DynamicsDisc (factored DARC)       & $-720$ \\
\bottomrule
\end{tabular}%
}
\end{table*}

\begin{table*}[t]
\centering
\small
\caption{Ablation on SIM online (Group C), mean cumulative reward over 4--5 seeds. SIM-trained policies evaluated on the SUMO target environment. Bold: best.}
\label{tab:ablation_sim}
\resizebox{\textwidth}{!}{%
\begin{tabular}{llr}
\toprule
Config & Components & Cum.\ reward (K) \\
\midrule
\texttt{sim\_is}            & Contrastive only                   & $\mathbf{-658}$ \\
\texttt{sim\_dyn\_t05}      & DynamicsDisc, temp$=0.5$           & $-669$ \\
\texttt{sim\_darc}          & TransDisc only                     & $-674$ \\
\texttt{sim\_is\_calql}     & Contrastive + Q-floor              & $-680$ \\
\texttt{sim\_dyn}           & DynamicsDisc, temp$=1.0$           & $-682$ \\
\texttt{sim\_baseline}      & no IS, no Q-floor                  & $-684$ \\
\texttt{sim\_calql}         & no IS, Q-floor only                & $-692$ \\
\texttt{sim\_darc\_calql}   & TransDisc + Q-floor                & $-694$ \\
\texttt{sim\_dyn\_t20}      & DynamicsDisc, temp$=2.0$           & $-699$ \\
\texttt{sim\_full}          & Contrastive + Q-floor + KL + warmup & $-705$ \\
\texttt{sim\_rescale}       & TransDisc + sim-reward rescale     & $-726$ \\
\bottomrule
\end{tabular}%
}
\end{table*}

\paragraph{SUMO-online vs SIM-online: training-environment matters}
The discriminator preference flips between training environments: TransitionDisc is best in SUMO-online (\Cref{tab:ablation}), Contrastive is best in SIM-online (\Cref{tab:ablation_sim}).
This is consistent with the H2O+ premise that the discriminator must distinguish the online (training) environment from the offline-data distribution; when the online environment differs more from offline (as in SIM-online, where the LSTM sim has different dynamics than the SUMO-collected offline data), a more flexible action-marginalised discriminator (Contrastive) helps; when online is closer to offline (SUMO-online uses the same simulator the offline data was collected from), the simpler TransitionDisc suffices and adding the more flexible variant slightly hurts.

\subsection{Analysis}
\label{sec:exp:analysis}

\Cref{fig:analysis} reports three diagnostic curves derived from the SUMO-online training logs (smoothed across seeds).

\paragraph{IS weight evolution (panel a)}
The realised IS weight $\wis = (\sqrt{w_{\text{IS}}})^2$ tracks the discriminator's binary-classifier ratio averaged across each minibatch, on a logarithmic scale.
TransitionDiscriminator stabilises around $0.2$--$0.3$ (heavily suppressing online samples), DynamicsDiscriminator around $0.4$--$0.6$, and the Contrastive InfoNCE variant around $0.7$--$0.9$.
This separation is consistent with the discriminator-choice ordering observed in \Cref{tab:ablation,tab:ablation_sim}: low IS weight (TransDisc) is appropriate when online and offline distributions are close (SUMO-online setting), while higher IS weight (Contrastive) is needed when they differ (SIM-online).

\paragraph{Discriminator convergence (panel b)}
All three discriminator variants converge in classification loss within the first $2{-}3$K gradient steps, but at different plateau values: TransDisc settles lowest (best classification), DynamicsDisc highest (worst classification, consistent with its weakest performance in \Cref{tab:ablation,tab:ablation_sim}), and Contrastive intermediate.
The Contrastive InfoNCE plateau, with bounded sensitivity to action via construction, is a stable training signal even when classification accuracy is moderate.

\paragraph{Training stability across H2O+ variants (panel c)}
The mean real-batch reward (rolling-averaged) tracks the policy's value during training across four representative variants.
Configurations including the Q-floor anchor ($\texttt{sumo\_calql}$, $\texttt{sumo\_full}$) maintain a more stable trajectory through the full $200$ epochs, while configurations without it ($\texttt{sumo\_baseline}$, $\texttt{sumo\_darc}$) exhibit small late-training drifts.
The drift magnitude is small relative to the cross-config gap reported in the SUMO-eval ablation (\Cref{tab:ablation}), explaining why the Q-floor's effect on final SUMO eval is configuration-dependent rather than uniformly positive.

\paragraph{Action invariance: learned-model diagnostic and data-level check}
We measure the dependence of the discriminator's output on the action via the ratio $f_\psi(s, a, s') / f_\psi(s, a_{\text{perm}}, s')$ with random action permutation $a_{\text{perm}}$ across the batch.
Contrastive achieves $0.96$ (action-invariant by construction; the score function does not see $a$); TransitionDiscriminator achieves $0.79$ (action-conditioned but with weak signal under transit's exogenous-dominated transitions).
This is a property of the \emph{learned model}; we additionally report a direct data-level test of \Cref{ass:action_invariance} that does not depend on a trained discriminator.
We discretise actions into $K_a = 3$ bins (no-hold, short-hold, long-hold), project $(s, s')$ to a $5 \times 5$ PCA grid, and within each $(s, s')$ cell that contains samples in multiple action bins compute the action-dependence ratio
$\hat{r} = \mathrm{mean}_{(s_{\text{in}}, s_{\text{out}})}\bigl[\,\mathrm{std}_a \log \hat\rho\,\bigr] / \mathrm{mean}_a\bigl[\,\mathrm{std}_{(s_{\text{in}}, s_{\text{out}})} \log \hat\rho\,\bigr]$
where $\hat\rho = \hat P_{\text{real}}(s' \mid s, a) / \hat P_{\text{sim}}(s' \mid s, a)$.
A small ratio indicates that variation across $a$ is small relative to variation across $(s, s')$, supporting \Cref{ass:action_invariance}.
On $50{,}000$ offline-side transitions paired with sim-core rollouts under matching initial-state seeds, we measure $\hat{r} = 0.156$ at the $5 \times 5$ grid (and $\hat{r} = 0.169$ at $6 \times 6$); we adopt $0.30$ as a working threshold (small ratio = action-variation dominated by $(s, s')$-variation), and treat the result as supportive evidence rather than conclusive verification.
Supporting heatmaps and per-cell ratios accompany the artefact release; the script \texttt{SimpleSAC/verify\_assumption1.py} reproduces the analysis end-to-end.
This data-level result is consistent with the contrastive variant's strong empirical performance under SIM-online but is not a guarantee on regions of state space underrepresented in $\Doff$.

\begin{figure*}[t]
\centering
\includegraphics[width=\linewidth]{figures/analysis.pdf}
\caption{Training-time diagnostics, computed from \emph{SUMO-online} training logs (the same-fidelity oracle setting; corresponding SIM-online curves are qualitatively similar in IS-weight ordering and discriminator-loss plateau but quantitatively differ since the SIM-vs-offline distribution gap is larger). (a)~Effective IS weight $\wis$ over training steps for the three discriminator variants (mean $\pm$ std across SUMO-online seeds; log scale). TransDisc converges to a low weight ($\approx 0.25$), DynDisc to a moderate weight, and Contrastive InfoNCE to a higher weight ($\approx 0.7$). (b)~Discriminator classification loss; lower is better classification of the offline/online split. (c)~Real-batch reward over training across four H2O+ variants; configurations with Q-floor anchor maintain greater late-training stability.}
\label{fig:analysis}
\end{figure*}


\subsection{Sensitivity}
\label{sec:exp:sensitivity}

\paragraph{Ensemble size}
Varying $\Ecnt \in \{3, 5, 8\}$ leaves the final reward essentially unchanged ($< 1\%$ range), though $\Ecnt{=}3$ shows higher per-seed variance.
Training time scales sub-linearly in $\Ecnt$ (independent target computation amortises across heads).

\paragraph{Data composition}
Group D of the ablation removes one offline source at a time.
Removing the SAC-trained source most hurts online performance (loss of optimal-action coverage); removing the weak-random source least hurts (coverage is redundant with heuristic); removing the RE-SAC expert ep39 source removes the Q-floor anchor's most-calibrated samples.
All three individual-source drops remain above the H2O+ baseline.

\paragraph{Data-efficiency curve}
\Cref{fig:data_efficiency_fig} visualises the relationship.

\begin{figure}[t]
\centering
\includegraphics[width=0.85\linewidth]{figures/data_efficiency.pdf}
\caption{Data efficiency: H2O+ vs pure offline RL on SUMO eval, varying offline buffer size from 100K to 675K. Bars are mean$\pm$std over 5 seeds. H2O+ at 200K matches pure offline at 675K on the \emph{single default scenario} used here (SUMO seed $42$, default OD), implying a $3.4\times$ data-efficiency benefit from the online phase under this protocol; this ratio is computed on the default-scenario protocol and is not measured on the $3 \times 3$ scenario grid of \Cref{tab:multiseed}. Pure online RL (red dashed) is included for reference.}
\label{fig:data_efficiency_fig}
\end{figure}

\Cref{tab:data_size} reports Group E: cumulative reward as a function of offline buffer size for both pure offline (Cal-QL) and the H2O+ SIM-online pipeline.
Two patterns emerge.
First, both curves plateau at $\approx 200$K offline transitions; quadrupling data to 400K--675K yields essentially no improvement under either method, suggesting the offline buffer's bus-trajectory diversity rather than its raw size is the bottleneck.
Second, H2O+ consistently improves over pure offline across all sizes by $45$--$92$K cumulative reward, and reaches with $200$K transitions a level ($-686$K) that pure offline cannot match even with $675$K transitions ($-750$K).
The online SIM phase therefore substitutes for offline-data volume in a quantitative sense (a $3.4\times$ data-efficiency gain at the 200K vs 675K comparison point).

\begin{table*}[t]
\centering
\small
\caption{Group E: Data-efficiency. Cumulative SUMO reward (K) vs offline buffer size, 5 seeds each. H2O+ reaches with 200K what pure offline never reaches with 675K. The H2O+ column uses the \texttt{sim\_is\_calql} configuration (Contrastive + Q-floor) for historical reasons --- this sweep predates the SIM-online ablation that identified the Contrastive-only variant as the deployment-best (\Cref{tab:ablation_sim}). The data-efficiency \emph{ratio} (H2O+ at 200K matching pure offline at 675K) is robust to recipe choice within the H2O+ family.}
\label{tab:data_size}
\resizebox{\textwidth}{!}{%
\begin{tabular}{lrrr}
\toprule
Offline size & Pure offline (Cal-QL) & H2O+ (Contrastive + Q-floor) & Gap \\
\midrule
100K & $-789 \pm 42$ & $-697 \pm 61$ & $+92$ \\
200K & $-750 \pm 14$ & $-686 \pm 48$ & $+64$ \\
400K & $-749 \pm 24$ & $-689 \pm 32$ & $+60$ \\
Full ($\approx$ 675K) & $-750 \pm 28$ & $-705 \pm 38$ & $+45$ \\
\bottomrule
\end{tabular}%
}
\end{table*}

\paragraph{KL anchor and warmup}
Enabling the KL anchor to a 50-epoch policy snapshot \citep{zhou2024wsrl} stabilises early training but has small final impact; disabling it yields faster but slightly noisier convergence.

\subsection{Multi-scenario robustness and operational metrics}
\label{sec:exp:multiseed}

The single-seed/single-OD numbers reported in \Cref{tab:main} are reproducible but narrow.
To address generalisation across passenger-demand days and SUMO seeds we additionally
evaluate each method across a $3 \times 3$ scenario grid: SUMO seeds
1001, 1002, and 1003 crossed with OD-scale multipliers 0.6, 0.8, and 1.0
(applied uniformly to both vehicle and passenger demand via SUMO's
\texttt{--scale} flag).
Seventeen method rows are compared on the same scenario set under common random numbers (\Cref{tab:multiseed}):
ep39 (the RE-SAC SUMO expert checkpoint included in $\Doff$);
two SUMO-online H2O+ checkpoints from \Cref{tab:ablation} (TransDisc-only and TransDisc + Q-floor; same-fidelity oracle);
the deployment-relevant SIM-online H2O+ Contrastive-only variant ($3$ training seeds);
five standard offline-to-online RL baselines trained with three seeds each (RLPD, WSRL, TD3+BC, IQL, AWAC);
two pure-online SAC checkpoints (a legacy $200$-epoch single-seed run, and a fresh $1{,}000$-epoch retrain across $3$ training seeds);
two behaviour-cloning variants (the full $\Doff$ at one training seed, and the ep39-only subset at three training seeds);
zero-hold (no intervention);
and Daganzo cooperative holding at three $\alpha$ values ($0.3$, $0.4$, $0.6$).
Each row reports a single underlying checkpoint (or code-based controller) evaluated on the same $9$ scenarios, except the seven multi-seed rows (SIM-online H2O+ Contrastive, $1{,}000$-epoch pure-online SAC, RLPD, WSRL, TD3+BC, IQL, AWAC, and BC ep39-only), each of which uses $3$ training-seed checkpoints $\times$ $9$ scenarios $= 27$ evaluations.
The total evaluation count is $9 \times 9 + 8 \times 27 = 297$.
The reported $\pm$ in \Cref{tab:multiseed} is the standard error \emph{across the $n$ paired evaluations of that row}; for the $n=27$ rows this includes both training-seed and scenario variation, while for the $n=9$ rows it includes only scenario variation (the underlying checkpoint or controller is fixed). We do not perform a hierarchical bootstrap that separates the two sources of variance; readers should treat the $\pm$ as a descriptive within-row spread, not as a uniformly inferential quantity, when comparing rows of different $n$.

\begin{table*}[t]
\centering
\scriptsize
\setlength{\tabcolsep}{3pt}
\caption{Multi-scenario evaluation: $3$ SUMO seeds $\{1001, 1002, 1003\} \times 3$ OD-scale multipliers $\{0.6, 0.8, 1.0\}$ $= 9$ scenarios per method (paired under common random numbers). Cumulative reward in K (mean $\pm$ standard error across the $n$ paired scenarios; $n = 9$ unless noted). Gain vs zero-hold is the paired mean improvement in cumulative reward over the zero-hold controller on the same scenario grid, also in K; this column makes the stress-grid reward scale comparable across methods. Passenger waiting is mean per-leg station waiting time; HW CV is the cross-route mean of per-line headway coefficient of variation; large-gap rate is the fraction of headways exceeding $1.5\times$ scheduled; Jain's index is computed over per-line headway-CV values (higher = more equitable across lines). Bold = best per column.}
\label{tab:multiseed}
\resizebox{\textwidth}{!}{%
\begin{tabular}{lrrrrrr}
\toprule
Method & Cum.\ reward (K) & Gain vs zero (K) & Pax wait (s) & HW CV & Large-gap & Jain \\
\midrule
\multicolumn{7}{l}{\emph{Expert reference policy (transitions in $\Doff$)}} \\
RE-SAC SUMO expert ep39 (ckpt in $\Doff$)                         & $\mathbf{-1{,}634 \pm 253}$ & $\mathbf{+593}$ & $475.0$ & $\mathbf{0.627}$ & $\mathbf{0.000}$ & $0.913$ \\
\midrule
\multicolumn{7}{l}{\emph{Deployment-relevant learned policy (Stage 2 in cheap simulator)}} \\
\textbf{H2O+ Contrastive only (SIM-online, $3$ seeds, $n{=}27$)}  & $-1{,}725 \pm 150$ & $+502$ & $499.8$ & $0.646$ & $0.001$ & $0.912$ \\
\midrule
\multicolumn{7}{l}{\emph{Same-fidelity oracle (Stage 2 in target simulator; not deployable)}} \\
H2O+ TransDisc only (SUMO-online oracle)                          & $-1{,}682 \pm 269$ & $+546$ & $494.8$ & $0.644$ & $\mathbf{0.000}$ & $0.907$ \\
H2O+ TransDisc + Q-floor (SUMO-online ablation)                   & $-1{,}697 \pm 272$ & $+530$ & $491.1$ & $0.640$ & $0.001$ & $0.912$ \\
\midrule
\multicolumn{7}{l}{\emph{Standard offline-to-online RL baselines ($3$ seeds each, $n{=}27$; snapshot-reset off, matching H2O+ Contrastive)}} \\
RLPD \citep{ball2023efficient} (offline+online $50/50$ mix)        & $-1{,}726 \pm 145$ & $+501$ & $506.9$ & $0.644$ & $0.001$ & $0.917$ \\
WSRL \citep{zhou2024wsrl} ($20$-ep warmup-collect)                  & $-1{,}709 \pm 147$ & $+518$ & $503.7$ & $0.647$ & $0.001$ & $0.915$ \\
TD3+BC \citep{fujimoto2021minimalist} (pure offline)               & $-1{,}746 \pm 153$ & $+481$ & $492.0$ & $0.656$ & $0.001$ & $0.913$ \\
IQL \citep{kostrikov2021offline} (Q + V expectile + AWR)           & $-2{,}105 \pm 110$ & $+122$ & $483.0$ & $0.662$ & $0.009$ & $0.919$ \\
AWAC \citep{nair2020awac} (Q + AWR with batch baseline)            & $-2{,}198 \pm 88$  & $+29$ & $482.3$ & $0.665$ & $0.012$ & $0.923$ \\
\midrule
\multicolumn{7}{l}{\emph{Other learned policies}} \\
Pure-online SAC ($1{,}000$ ep retrain, $3$ seeds, $n{=}27$)       & $-1{,}739 \pm 152$ & $+488$ & $505.0$ & $0.650$ & $0.001$ & $0.912$ \\
Pure-online SAC ($200$ ep, single seed, legacy)                   & $-2{,}019 \pm 166$ & $+208$ & $\mathbf{447.1}$ & $0.628$ & $0.016$ & $0.927$ \\
BC (ep39-only, $3$ seeds, $n{=}27$)                                & $-1{,}659 \pm 142$ & $+568$ & $482.0$ & $0.634$ & $\mathbf{0.000}$ & $0.913$ \\
BC (full $\Doff$, $1$ training seed)                              & $-2{,}244 \pm 174$ & $-16$ & $507.0$ & $0.693$ & $0.010$ & $0.921$ \\
\midrule
\multicolumn{7}{l}{\emph{Unlearned baselines}} \\
Zero-hold                                                          & $-2{,}227 \pm 151$ & $0$ & $485.3$ & $0.673$ & $0.013$ & $0.926$ \\
Daganzo cooperative ($\alpha = 0.3$, retuned)                     & $-2{,}436 \pm 175$ & $-209$ & $488.1$ & $0.766$ & $0.011$ & $0.936$ \\
Daganzo cooperative ($\alpha = 0.4$, retuned)                     & $-2{,}469 \pm 179$ & $-242$ & $487.8$ & $0.774$ & $0.012$ & $0.936$ \\
Daganzo cooperative ($\alpha = 0.6$, textbook)                    & $-2{,}511 \pm 167$ & $-284$ & $494.0$ & $0.781$ & $0.010$ & $\mathbf{0.938}$ \\
\bottomrule
\end{tabular}%
}
\end{table*}

The cumulative-reward magnitudes in \Cref{tab:multiseed} are systematically more negative than the corresponding rows in \Cref{tab:main} ($\approx -1{,}700$K vs $\approx -650$K for H2O+).
This is expected: \Cref{tab:main} reports a single SUMO seed under the default OD profile, while \Cref{tab:multiseed} averages $9$ scenarios that include lower OD scales ($0.6$, $0.8$) where headway penalties dominate the composite reward and seeds $\{1001, 1002, 1003\}$ that were not used during training.
For this reason \Cref{tab:multiseed} also reports a paired gain over zero-hold: the H2O+ rows still improve by $+502$K to $+546$K on the stress grid, so the large negative raw-return values should not be read as ``no control'' behaviour.
The bottom-of-table ranking is preserved (zero-hold, BC, and Daganzo always trail by a wide margin), but the top-of-table ranking changes between protocols: \Cref{tab:main} has H2O+ narrowly above ep39 (single scenario), while \Cref{tab:multiseed} has ep39 narrowly above H2O+ on the scenario-grid mean (with the H2O+ vs ep39 ordering depending on OD scale, see below).
Four observations from \Cref{tab:multiseed} sharpen the headline claim:

\paragraph{Cumulative reward: ep39, H2O+, and $1{,}000$-epoch pure-online SAC form a near-tie at the top.}
Across the $9$ paired scenarios under common random numbers, the four strongest methods cluster within $\approx 100$K of one another:
ep39 (the RE-SAC SUMO expert checkpoint included in $\Doff$) is nominally first at $-1{,}634 \pm 253$K, ahead of H2O+ TransDisc-only ($-1{,}682 \pm 269$K, paired $\Delta = -47$K vs ep39, $95\%$ bootstrap CI $[-77, -12]$ excludes zero, although the Wilcoxon signed-rank $p = 0.055$ is at the borderline of conventional significance under $n = 9$; we read the CI as the more informative summary here);
H2O+ TransDisc + Q-floor sits between them ($-1{,}697 \pm 272$K, paired $\Delta = +16$K above H2O+ TransDisc-only, CI $[+5, +27]$, $p = 0.027$---a small but statistically significant gap);
$1{,}000$-epoch pure-online SAC averaged over $3$ training seeds reaches $-1{,}739 \pm 152$K (paired $\Delta = +57$K vs SUMO-online H2O+ TransDisc-only, $p = 0.004$);
the deployment-relevant SIM-online H2O+ Contrastive-only variant (averaged over $3$ training seeds, $n = 27$) reaches $-1{,}725 \pm 150$K, paired $\Delta = +43$K below the SUMO-online oracle ($95\%$ CI $[+29, +58]$, $p < 0.001$) --- this $43$K margin quantifies the cost of training in the cheap simulator instead of in the target.
H2O+ does \emph{not} beat ep39 in this multi-scenario protocol; the headline ``slightly above ep39'' from the single-scenario \Cref{tab:main} ($-658$K vs $-666$K) does not extrapolate to the $3 \times 3$ stress grid.

The headline gap H2O+ \emph{does} produce reliably is over the unlearned baselines:
$25\%$ over zero-hold ($-2{,}227$K) and BC ($-2{,}244$K);
$31$--$33\%$ over Daganzo at any of the three tested $\alpha$ values ($-2{,}436$K to $-2{,}511$K).
Cross-method paired Wilcoxon tests (\Cref{tab:paired_stats}) confirm: H2O+ TransDisc-only is significantly better than every non-ep39 comparator (all $p \leq 0.027$); against ep39 the comparison is borderline-significant in the wrong direction.

Because ep39's transitions are part of $\Doff$ (\Cref{sec:exp:setup}), this finding is best read as: \emph{H2O+ recovers most of the RE-SAC SUMO expert's performance and reliably beats the unlearned controllers, but does not consistently surpass that expert under multi-scenario evaluation}.
H2O+'s value here is therefore that it transfers and fine-tunes a strong expert-derived offline buffer under a tunable optimisation target, not that it discovers strictly better holding behaviour than the RE-SAC expert across the scenario set we tested.

\begin{table*}[t]
\centering
\small
\caption{Paired-scenario statistics with the \emph{same-fidelity oracle} (SUMO-online H2O+ TransDisc-only) as base, on the $9$ paired scenarios of \Cref{tab:multiseed}. Mean difference reported with bootstrap $95\%$ CI ($2{,}000$ resamples) and Wilcoxon signed-rank $p$-value. Positive mean difference = the SUMO-online oracle H2O+ has higher (better) cumulative reward than the comparator. This table is informative for ranking H2O+ variants under matched-fidelity training but does not directly answer the deployment question; \Cref{tab:deployment_paired} reports the same comparisons with the deployment-relevant SIM-online H2O+ Contrastive variant as base.}
\label{tab:paired_stats}
\resizebox{\textwidth}{!}{%
\begin{tabular}{lrrr}
\toprule
Comparator & $\Delta$ mean (K) & 95\% CI & Wilcoxon $p$ \\
\midrule
Daganzo ($\alpha = 0.6$, textbook) & $+829$ & $[+656,\ +1{,}046]$ & $0.004$ \\
BC (full $\Doff$)                  & $+562$ & $[+395,\ +754]$     & $0.004$ \\
Zero-hold                          & $+546$ & $[+343,\ +786]$     & $0.004$ \\
Pure-online SAC (legacy $200$ ep)  & $+337$ & $[+157,\ +541]$     & $0.004$ \\
Pure-online SAC ($1{,}000$ ep retrain, $3$ training seeds) & $+57$ & $[+29,\ +88]$ & $0.004$ \\
H2O+ Contrastive only (SIM-online, $3$ training seeds, $n{=}27$) & $+43$ & $[+29,\ +58]$ & $<0.001$ \\
H2O+ TransDisc + Q-floor           & $+16$  & $[+5,\ +27]$        & $0.027$ \\
RE-SAC SUMO expert ep39                 & $-47$ & $[-77,\ -12]$  & $0.055$ \\
\bottomrule
\end{tabular}%
}
\end{table*}

\paragraph{Deployment-focused statistics: SIM-online H2O+ Contrastive as base}
The comparison most relevant to a real deployment is not against the SUMO-online oracle (which assumes Stage~2 can train in the same high-fidelity simulator used for evaluation) but against the SIM-online H2O+ Contrastive variant, which is the only learned configuration in our matrix that an operator could realistically retrain inside the cheap LSTM sim-core $\Msim$ on a daily basis.
\Cref{tab:deployment_paired} reports paired Wilcoxon comparisons with this deployment-relevant policy as the base, on the $9$ paired scenarios.
Three patterns are worth separating.
First, against the RE-SAC SUMO expert ep39, SIM-online H2O+ loses by a paired mean of $-90$K ($95\%$ CI $[-117, -65]$, $p < 0.001$): under realistic training-budget constraints H2O+ does not reach the high-fidelity expert checkpoint, and this gap is statistically robust.
Second, against the same-fidelity SUMO-online oracle, SIM-online H2O+ pays a paired mean of $-43$K ($95\%$ CI $[-58, -29]$, $p < 0.001$): roughly half of the gap-to-ep39 is explainable as a fidelity penalty for training in $\Msim$ rather than $\Mreal$.
Third, against the $1{,}000$-epoch pure-online SAC retrained from scratch in $\Msim$, SIM-online H2O+ has a small positive paired mean of $+14$K ($95\%$ CI $[+4, +25]$, $p = 0.021$); the CI excludes zero and the test rejects equality at the conventional threshold, so the offline-pretrained variant has a small but statistically significant edge over the from-scratch SAC trained for the same wall-clock budget.
Fourth, against all unlearned baselines (zero-hold, BC, three Daganzo variants), SIM-online H2O+ wins by paired means of $+503$ to $+786$K, all at $p < 0.001$---the deployment-relevant H2O+ variant is reliably and substantially better than any controller that does not use both offline data and online interaction.

\begin{table*}[t]
\centering
\small
\caption{Deployment-focused paired-scenario statistics: SIM-online H2O+ Contrastive (the only learned configuration in our matrix that retrains inside the cheap LSTM sim-core, $3$ training seeds, $n=27$ paired evaluations) vs each comparator on the same-scenario subset where both methods have data. Bootstrap $95\%$ CI ($2{,}000$ resamples) and Wilcoxon signed-rank $p$-value. Positive mean difference = SIM-online H2O+ has higher (better) cumulative reward. The first three rows are the comparisons most relevant for a deployment decision; the bottom rows are sanity-check margins over unlearned controllers.}
\label{tab:deployment_paired}
\resizebox{\textwidth}{!}{%
\begin{tabular}{lrrr}
\toprule
Comparator & $\Delta$ mean (K) & 95\% CI & Wilcoxon $p$ \\
\midrule
\multicolumn{4}{l}{\emph{Vs.\ reference and oracle}} \\
RE-SAC SUMO expert ep39                             & $-90$  & $[-117,\ -65]$ & $<0.001$ \\
BC (ep39-only, $3$ seeds)                            & $-65$  & $[-90,\ -39]$  & $<0.001$ \\
H2O+ TransDisc + Q-floor (SUMO-online ablation)       & $-27$  & $[-38,\ -16]$  & $<0.001$ \\
H2O+ TransDisc-only (SUMO-online same-fidelity oracle) & $-43$  & $[-58,\ -29]$  & $<0.001$ \\
\midrule
\multicolumn{4}{l}{\emph{Vs.\ standard offline-to-online RL baselines (snapshot-reset off)}} \\
WSRL \citep{zhou2024wsrl}                            & $-15$  & $[-27,\ -2]$    & $0.021$ \\
RLPD \citep{ball2023efficient}                       & $+2$   & $[-13,\ +17]$   & $0.991$ \\
TD3+BC \citep{fujimoto2021minimalist}                & $+21$  & $[+5,\ +38]$    & $0.036$ \\
Pure-online SAC ($1{,}000$ ep retrain, $3$ seeds)    & $+14$  & $[+4,\ +25]$    & $0.021$ \\
IQL \citep{kostrikov2021offline} (pure offline)      & $+380$ & $[+274,\ +488]$ & $<0.001$ \\
AWAC \citep{nair2020awac} (pure offline)             & $+474$ & $[+349,\ +598]$ & $<0.001$ \\
\midrule
\multicolumn{4}{l}{\emph{Vs.\ unlearned baselines}} \\
Zero-hold                                            & $+503$ & $[+380,\ +633]$ & $<0.001$ \\
BC (full $\Doff$, $1$ seed)                           & $+519$ & $[+420,\ +623]$ & $<0.001$ \\
Daganzo ($\alpha = 0.3$, retuned)                    & $+712$ & $[+607,\ +819]$ & $<0.001$ \\
Daganzo ($\alpha = 0.4$, retuned)                    & $+745$ & $[+643,\ +854]$ & $<0.001$ \\
Daganzo ($\alpha = 0.6$, textbook)                   & $+786$ & $[+672,\ +903]$ & $<0.001$ \\
\bottomrule
\end{tabular}%
}
\end{table*}

The deployment-focused reading of these two tables is: an operator cannot choose the SUMO-online oracle (it requires training inside the high-fidelity simulator that is also the evaluation environment), so the relevant policy on offer from this work is SIM-online H2O+ Contrastive. To keep the comparison configuration-clean, RLPD and WSRL are evaluated with \texttt{snapshot\_reset} switched off (matching H2O+ Contrastive); the snapshot-reset choice itself perturbs RLPD/WSRL by only $\approx 1$--$2$K paired mean, well inside per-seed noise (see Robustness paragraph below).

\paragraph{Standard offline-to-online RL baselines: H2O+ is competitive but not uniformly best}
The standard offline-to-online RL comparison includes IQL, AWAC, TD3+BC (all offline-only), and RLPD, WSRL (offline+online hybrid). All run with three training seeds and the same $9$-scenario evaluation grid; each is implemented to operate on the same $675$K-transition $\Doff$ buffer. The picture for H2O+ Contrastive (SIM-online) on cumulative reward is:
(i) it \emph{dominates} the pure-offline baselines IQL ($+380$K, $p < 0.001$), AWAC ($+474$K, $p < 0.001$), and BC on the full mixed buffer ($+519$K, $p < 0.001$), confirming that the online phase is essential for the bus environment regardless of which offline-to-online algorithm is used;
(ii) it has a small statistically significant edge over TD3+BC ($+21$K, $p = 0.036$) and $1{,}000$-epoch pure-online SAC ($+14$K, $p = 0.021$);
(iii) it is \emph{indistinguishable} from RLPD ($+2$K, $p = 0.99$) and is \emph{slightly outperformed} by WSRL ($-15$K, $p = 0.021$).
The (iii) finding is the most informative: RLPD and WSRL are both H2O+ with the discriminator, IS reweighting, and offline-Q floor switched off. That neither falls behind H2O+ on this benchmark suggests the explicit dynamics correction does not provide a measurable additional benefit on top of a hybrid offline+online buffer with a $50/50$ mix and standard SAC updates --- the IS weights converge to near-uniform values and the discriminator signal is small relative to within-batch variance (\Cref{tab:ablation,fig:analysis}). H2O+ remains a sensible recipe (it never substantially underperforms), but the discriminator-driven correction is not the source of competitive advantage on this transit benchmark. We treat this as a useful negative result rather than a recipe contradiction; \Cref{sec:discussion} discusses what it implies for transferring H2O+ to other multi-fidelity domains.
For these RLPD/WSRL numbers we ran an extra robustness experiment: we trained each baseline a second time with \texttt{snapshot\_reset} \emph{off}, matching the H2O+ Contrastive configuration (\texttt{snapshot\_reset} resets the cheap-simulator episode initial state from a SUMO snapshot with probability $0.5$). The paired-difference between the two configurations is $+2.3$K for RLPD ($p = 0.31$) and $+1.1$K for WSRL ($p = 0.56$), well inside per-seed noise, so the comparisons reported above are robust to the snapshot-reset choice.

A second observation worth flagging is that BC on the ep39-only subset ($-1{,}659 \pm 142$K) reaches paired difference of only $-25$K against ep39 itself, much closer than BC on the full mixed buffer ($-2{,}244$K, $-610$K below ep39). This suggests that mixing ep39 transitions with $\epsilon$-random and weak-heuristic transitions in $\Doff$ \emph{degrades} BC --- the imitation target is averaged across heterogeneous behaviour qualities. None of the offline-to-online RL methods evaluated here (H2O+, RLPD, WSRL, TD3+BC) face this BC-style degradation because they reweight via Q values rather than imitate uniformly.

\paragraph{Operational metrics: ep39 best on headway regularity, legacy $200$-epoch SAC best on passenger waiting}
The cumulative-reward ranking does not transfer cleanly to the passenger-side columns.
The legacy $200$-epoch single-seed pure-online SAC checkpoint has the lowest mean passenger waiting time across the $9$ scenarios ($447.1$s, vs $499.1$s for SIM-online H2O+ Contrastive and $475.0$s for ep39); ep39 has the lowest cross-route headway CV ($0.627$, vs $\approx 0.64$--$0.65$ for the H2O+ and $1{,}000$-epoch SAC variants); the $1{,}000$-epoch retrain of pure-online SAC sits at $505.0$s waiting and $0.650$ HW CV, slightly worse than its $200$-epoch predecessor on these passenger-side metrics despite being substantially stronger on the composite training reward.
This non-monotonicity --- training longer on the shaped reward improves cumulative reward but does \emph{not} monotonically improve passenger waiting or headway CV --- is informative on its own: the composite reward used during training penalises holding aggressively (per-second hold cost dominates the bunching reduction the hold buys); when we examine the hold-time distributions across methods (means $37$--$45$s, $90$th percentiles $52$--$60$s), the four top RL methods are statistically similar in \emph{how long} they hold, so the passenger-wait differences are driven by \emph{when} and \emph{which buses} are held rather than gross holding intensity.
The implication is a calibration gap, not a methodological failure: \emph{the optimisation objective the paper trains for is not strictly aligned with per-leg passenger waiting, and at $1{,}000$-epoch budgets pure-online SAC and H2O+ over-fit the shaped reward at modest cost on waiting}.
A reward re-weighting that aligns with per-leg waiting cost would change the optimisation landscape across the entire ablation matrix; we flag this as a calibration step required before any operator pilot (\Cref{sec:disc:operational}) and visualise the trade-off in \Cref{fig:pareto}: ep39 sits at moderate waiting / high reward, the H2O+ rows cluster nearby, the legacy $200$-epoch pure-online SAC sits at low waiting / low reward, the $1{,}000$-epoch retrain at moderate waiting / mid reward, and Daganzo / BC / zero-hold occupy the high-waiting / low-reward corner.

\begin{figure*}[t]
\centering
\includegraphics[width=0.95\linewidth]{figures/pareto.pdf}
\caption{Reward vs operational utility: per-method means across $9$ scenarios (3 SUMO seeds $\times$ 3 OD scales). Lower passenger waiting and higher cumulative reward are both desirable; no single method dominates both axes. The RE-SAC SUMO expert ep39 has the best cumulative reward; the legacy $200$-epoch single-seed pure-online SAC has the lowest per-leg waiting ($447$s) but at substantially worse cumulative reward; the $1{,}000$-epoch retrain of pure-online SAC sits at moderate waiting ($505$s) and mid reward; H2O+ Contrastive (SIM-online) and the SUMO-online H2O+ variants cluster near ep39. The reward-vs-operational trade-off motivates the calibration discussion in \Cref{sec:disc:operational}.}
\label{fig:pareto}
\end{figure*}

\paragraph{Diagnostic passenger-side breakdown: reward alignment caveat}
The cumulative reward and the mean per-leg waiting time reported in \Cref{tab:multiseed} together cover only one slice of the passenger-side picture.
\Cref{tab:passenger_full} reports four diagnostic metrics aggregated over the same $9$-scenario grid (and over $3$ training seeds for the rows where $n=27$): completed passenger trips per $18{,}000$\,s episode, $90$th-percentile per-leg waiting time, mean in-vehicle delay (time spent on a bus, including any delay attributable to holding), and mean total passenger travel time (waiting + in-vehicle).
Three observations are immediate.
First, completed-trips counts are essentially identical across all eleven methods ($\approx 312$ trips/episode); none of the controllers shed demand under any tested OD scale, so passenger-throughput is not a useful discriminator on this benchmark.
Second, the RL methods that win on cumulative reward (ep39, H2O+ TransDisc-only, H2O+ TransDisc + Q-floor, $1{,}000$-epoch pure-online SAC, SIM-online H2O+ Contrastive) cluster between $1{,}180$ and $1{,}270$\,s of mean in-vehicle delay --- substantially higher than zero-hold ($847$\,s), BC ($865$\,s), and Daganzo ($896$--$904$\,s).
The hold actions that the shaped reward rewards (suppressing bunching, smoothing inter-line headway variance) directly extend the time each rider spends on a moving bus.
Third, when waiting and in-vehicle are summed, the RL methods incur \emph{larger} bridge-derived mean total passenger travel time ($\approx 1{,}700$\,s) than zero-hold ($1{,}333$\,s), BC ($1{,}372$\,s), and Daganzo ($1{,}384$--$1{,}398$\,s) --- a $\approx 25\%$ increase.
We treat this as a reward-alignment warning rather than a passenger-benefit claim: under the current SUMO passenger bridge, no RL method trained on the shaped reward dominates the unlearned baselines on total passenger travel time.
The result bounds the operational claim we make: H2O+'s value on this benchmark is in cumulative reward and headway regularity (Jain index, HW CV), not in directly minimising total passenger travel time, and a calibrated reward that internalises in-vehicle cost would be a prerequisite to any deployment recommendation that targets passenger travel time as the operational objective (\Cref{sec:disc:operational}).

\begin{table*}[t]
\centering
\small
\caption{Diagnostic passenger-side metrics on the $3\times 3$ scenario grid. \emph{Completed trips} is the number of passenger trips finished within the $18{,}000$\,s episode; \emph{wait p90} is the $90$th-percentile per-leg waiting time; \emph{in-vehicle} is mean time on bus (including delay attributable to holding); \emph{total travel time} is waiting + in-vehicle under the current SUMO passenger bridge. All times in seconds, mean across paired scenarios ($n$ as in \Cref{tab:multiseed}). Daganzo's excess-wait reference is undefined (the analytical rule does not consume a target headway in this implementation), so excess-wait is not reported. Bold = best per column among RL methods (top); separately bold among unlearned baselines (bottom) where it differs from the global best.}
\label{tab:passenger_full}
\resizebox{\textwidth}{!}{%
\begin{tabular}{lrrrrr}
\toprule
Method & $n$ & Completed trips & Wait p90 (s) & In-vehicle (s) & Total travel (s) \\
\midrule
\multicolumn{6}{l}{\emph{Reference / hybrid offline+online learned policies}} \\
RE-SAC SUMO expert ep39                             & $9$  & $\mathbf{312}$ & $1{,}034.7$ & $1{,}224.5$ & $1{,}699$ \\
H2O+ Contrastive (SIM-online, $3$ seeds)             & $27$ & $312$          & $1{,}094.3$ & $1{,}222.7$ & $1{,}723$ \\
H2O+ TransDisc-only (SUMO-online oracle)             & $9$  & $312$          & $1{,}090.1$ & $1{,}222.4$ & $1{,}717$ \\
H2O+ TransDisc + Q-floor (SUMO-online ablation)      & $9$  & $312$          & $1{,}068.0$ & $1{,}180.8$ & $1{,}672$ \\
RLPD (3 seeds)                                       & $27$ & $312$          & $1{,}108.6$ & $1{,}234.0$ & $1{,}739$ \\
WSRL (3 seeds)                                       & $27$ & $312$          & $1{,}092.6$ & $1{,}241.6$ & $1{,}741$ \\
Pure-online SAC ($1{,}000$ ep, $3$ seeds)            & $27$ & $312$          & $1{,}115.0$ & $1{,}219.6$ & $1{,}725$ \\
Pure-online SAC ($200$ ep, single seed, legacy)      & $9$  & $312$          & $\mathbf{969.9}$ & $1{,}270.5$ & $1{,}718$ \\
\midrule
\multicolumn{6}{l}{\emph{Pure-offline policies (low in-vehicle, low total travel---they hold less)}} \\
TD3+BC (3 seeds)                                     & $27$ & $312$ & $1{,}089.4$ & $\mathbf{1{,}216.0}$ & $\mathbf{1{,}708}$ \\
BC (ep39-only, $3$ seeds)                            & $27$ & $312$ & $1{,}076.5$ & $1{,}240.4$ & $1{,}722$ \\
IQL (3 seeds)                                        & $27$ & $312$ & $1{,}018.8$ & $876.4$ & $1{,}359$ \\
AWAC (3 seeds)                                       & $27$ & $312$ & $1{,}008.5$ & $839.2$ & $1{,}321$ \\
\midrule
\multicolumn{6}{l}{\emph{Unlearned baselines (lower in-vehicle, lower total travel)}} \\
Zero-hold                                            & $9$  & $312$ & $958.6$ & $\mathbf{847.2}$ & $\mathbf{1{,}333}$ \\
BC (full $\Doff$, $1$ seed)                           & $9$  & $312$ & $1{,}021.7$ & $865.4$ & $1{,}372$ \\
Daganzo cooperative ($\alpha = 0.3$, retuned)        & $9$  & $312$ & $976.9$ & $896.4$ & $1{,}384$ \\
Daganzo cooperative ($\alpha = 0.4$, retuned)        & $9$  & $312$ & $972.8$ & $900.4$ & $1{,}388$ \\
Daganzo cooperative ($\alpha = 0.6$, textbook)       & $9$  & $312$ & $991.7$ & $904.0$ & $1{,}398$ \\
\bottomrule
\end{tabular}%
}
\end{table*}

\paragraph{Daganzo over-holds at the textbook $\alpha$}
Daganzo with the textbook $\alpha = 0.6$ achieves the highest Jain fairness ($0.938$, treating all lines symmetrically by design) but the worst cumulative reward, $\approx 250$K below zero-hold.
A per-bus diagnostic confirmed the policy fires occasional $60$-second clip-level holds on buses sandwiched between close neighbours; these are textbook-correct but the per-second holding penalty in our reward dominates the future bunching reduction.
Retuned $\alpha = 0.3$ on the same $9$-scenario grid reaches $-2{,}436 \pm 175$K and $\alpha = 0.4$ reaches $-2{,}469 \pm 179$K, supporting the interpretation as $\alpha$-tuning sensitivity rather than a structural failure of the analytical rule; the retuned variants remain below H2O+ but close most of the gap to zero-hold ($-2{,}227$K).

\paragraph{Pure-online SAC: training-budget caveat}
\Cref{tab:multiseed} reports both the legacy $200$-epoch pure-online SAC checkpoint used for parity comparison in \Cref{tab:main} ($-2{,}019$K, single training seed) and a fresh $1{,}000$-epoch retrain across $3$ training seeds ($-1{,}739$K, $n = 27$ scenario evaluations).
The $5\times$ longer training closes most of the gap to H2O+: the paired mean difference shrinks from $+337$K (at $200$ epochs) to $+57$K (at $1{,}000$ epochs), and pure-online SAC at $1{,}000$ epochs ranks third overall, ahead of zero-hold, BC, and all Daganzo variants.
The catastrophic-failure framing of \Cref{tab:main} therefore applies specifically to the $200$-epoch parity budget; ``pure-online SAC fails on this benchmark'' is a budget-conditioned statement, not a structural one.
At the $1{,}000$-epoch budget, pure-online SAC is a credible RL baseline that H2O+ beats only by a small (but statistically significant) margin on the shaped composite reward.

A per-scenario heatmap of cumulative reward (\Cref{fig:per_scenario_heatmap}) shows that the magnitude difference between \Cref{tab:main} ($\approx -650$K, default scenario) and \Cref{tab:multiseed} ($\approx -1{,}700$K, scenario-grid mean) is dominated by the OD $\in \{0.6, 0.8\}$ columns---under reduced demand the headway-deviation penalties remain in absolute units while fewer passengers ride per bus, so per-decision reward magnitudes inflate.
Bottom-of-ranking (BC, zero-hold, Daganzo always last three) is preserved cell by cell, but the top of the ranking shifts with OD:
\textbf{ep39 is best at OD $\in \{0.6, 0.8\}$ ($6/9$ cells)} and the two H2O+ variants share the top at OD $= 1.0$ ($3/9$ cells).
Pure-online SAC ($1{,}000$ ep) is consistently $4$th among the top RL methods.
The H2O+ vs ep39 ordering is therefore demand-dependent: H2O+ wins under high demand where active holding pays off, while ep39's lighter-touch policy is preferable under low-to-moderate demand where holding incurs net cost on the shaped reward.

\begin{figure}[t]
\centering
\includegraphics[width=0.95\linewidth]{figures/per_scenario_heatmap.pdf}
\caption{Per-scenario cumulative reward (K) on the $3 \times 3$ scenario grid. The OD = $0.6$ column produces the most-negative magnitudes for all methods, explaining the inflation between \Cref{tab:main} and \Cref{tab:multiseed}. The bottom of the ranking (BC, zero-hold, Daganzo $\alpha = 0.6$) is preserved cell by cell, but the top of the ranking shifts with OD: ep39 leads at OD $\in \{0.6, 0.8\}$ while the H2O+ variants share the top at OD $= 1.0$.}
\label{fig:per_scenario_heatmap}
\end{figure}

The scenario grid is a $3 \times 3$ stress test rather than an exhaustive multi-day simulation; remaining evaluation-scope gaps (weekend OD, incident days, train/calibration/test day splits) are discussed in \Cref{sec:disc:calibration}.

\subsection{Computational cost}
\label{sec:exp:compute}

Training one SIM-online configuration takes roughly $15$ minutes on a 2-core CPU allocation; one SUMO-online configuration takes roughly $3$ hours due to SUMO's libsumo single-session constraint.
The full sweep (approximately 30 distinct configurations $\times$ 4--5 seeds $\approx$ 150 runs) fits in approximately $120$ CPU-hours when parallelised across SIM/offline runs and serialised across SUMO runs, using the hardware-aware auto-parallel launcher (\Cref{app:impl}).
Evaluation of a single checkpoint on a full SUMO episode takes $\sim 2$ minutes.
